\subsection{Two Metric Space Theorems}
r we rly doing banach fixed pt lol.
hm. weird defn of a contraction. oh i lied its standard
\begin{defn}
	$f:M\gto M$ a \textbf{contraction} if $\ex\alpha\in[0,1)$ s.t.
	$d(f(x),f(y))\leq\alpha d(x,y)$ on all $M$.
\end{defn}

statement of banach fixed pt for completion
\begin{thm}[title=Banach Fixed-Point Theorem]
	sps $M$ complete, $f:M\gto M$ a contraction.
	then $\uex x\in M$ s.t. $f(x)=x$.
\end{thm} \

\begin{xmp}[source=Primary Source Material]
	let $(a_{ij})$ be doubly indexed real nums s.t.
	$\sup_{i,j\geq1}\abs{a_{ij}}\leq K$ and $\sum_{i,j}\abs{a_{ij}}<\infty$
	for some $K>0$. let $b\in\ell^{1}$ and consider:
	\begin{equation*}
		F:\ell^{1}\gto \ell^{1} \qquad (F(x))_{i}=\sum_{j}a_{ij}x_{j}^{2}+b_{j}
	\end{equation*}
	then:
	\begin{equation*}
		\sum_{i}(F(x))_{i}
		\ \leq \ \sum_{i}\left(\sum_{j}\abs{a_{ij}}\abs{x_{j}}^{2}+\abs{b_{i}}\right)
		\ \leq \ \norm{x}_{\infty}^{2}\sum_{ij}\abs{a_{ij}}+\abs{b}_{1}<\infty
	\end{equation*}
	we have $F:\ell^{1}\gto\ell^{1}$ is well-defn'd.
	we show that $\ex\ep>0$ s.t. if $\norm{b}_{1}\leq\ep$, then
	$\ex x\in\ell^{1}$ s.t. $x=F(x)$.

	idea: wtfind closed set $\subseteq\ell^{1}$ on which $F$ contracts.
	let $x,y\in\ell^{1}$. then:
	\begin{align*}
		\norm{F(x)-F(y)}_{1}
		& \ \leq \ \sum_{i}\abs{\sum_{j}a_{ij}(x_{j}^{2}-y_{j}^{2})} \\
		& \ \leq \ \sum_{i}\sum_{j}\abs{a_{ij}}\left(\abs{x_{j}}+\abs{y_{j}}\right)\abs{x_{j}-y_{j}} \\
		& \ \leq \ \norm{x-y}_{\infty}(\norm{x}_{\infty}+\norm{y}_{\infty})\left(\sum_{i,j}\abs{a_{ij}}\right) \\
		& \ \leq \ \norm{x-y}_{1}(\norm{x}_{1}+\norm{y}_{1})M
	\end{align*}
	note that as long as $x,y\in \bar{B}(0,r)$ with $r=1/(3M)$, then $F$ is a contraction.
	lastly, need to check that $F(\bar{B}(0,r))\subseteq\bar{B}(0,r)$.
	\begin{align*}
		\norm{F(x)}_{1}
		& \ \leq \ \sum_{i}\abs{\sum_{j}a_{ij}x^{2}+b_{i}} \\
		& \ \leq \ \norm{x}_{\infty}^{2}\sum_{i,j}\abs{a_{ij}}+\norm{b}_{1} \\
		& \ \leq \ \norm{x}_{1}^{2}M + \norm{b}_{1} \\
		& \ \leq \ \frac{M}{9M^{2}}+\ep
	\end{align*}
	whered that last bit come from. what.
	anyway as long as $M\geq1$ and $\ep\leq1/(9M)$:
	\begin{equation*}
		\norm{F(x)}_{1} \ \leq \ \frac{1}{9M}+\frac{1}{9M} \ \leq \ \frac{1}{3M}
	\end{equation*}
	so we can apply banach fixed pt.
\end{xmp}
point of the example: given non-linear fn's on a general space(?), it
may not be a contraction on the whole space.
however, we might be able to find a suff. small closed ball on which it is
in fact a contraction; if the fn also maps the ball to itself, then we're good.


\lecdate{Lec 12 - Feb 24 (Week 7)}

motivation: what does it mean for a pt in a space to be ``generic"?
for instance, in $\bR$, most pts are in some sense, probably irrational.
note:
\begin{equation*}
	\bar{\bQ}=\bR \qquad \bar{\bQ^{c}}=\bR \qquad \bar{\bQ\cap\bQ^{c}}=\eset
\end{equation*}
so density alone isnt a good test.

\begin{lm}
	let $A,B\subseteq M$ be open dense. then $A\cap B$ is open dense.
\end{lm} \

\begin{pf}[source=Primary Source Material]
	fix $x\in M$ and $\ep>0$.
	since $A$ dense, $\ex a\in A$ s.t. $a\in B(x,\ep)$.
	since $B(x,\ep)\cap A$ open, $\ex r>0$ s.t. $B(a,r)\subseteq B(x,\ep)\cap A$.
	since $B$ dense, $\ex b\in B(a,r)$.
	then $b\in A\cap B\cap B(x,\ep)$ as needed.
\end{pf}
we can now state bct.

\begin{thm}[title=Baire Category Theorem]
	sps $M$ complete. let $G_{n}$ be a seq of open dense subsets of $M$. then:
	\begin{equation*}
		G \ = \ \bigcap_{n=1}^{\infty}G_{n}
	\end{equation*}
	is dense.
\end{thm}
sets of this form have many names, such as:
\begin{itemize}
	\item thick.
	\item dense $G_{\delta}$ sets, i.e. a ctbl intersection of open sets
\end{itemize}

\begin{crll}
	if $M$ complete and:
	\begin{equation*}
		M=\bigcup_{n=1}^{\infty}F_{n}
	\end{equation*}
	where $F_{n}$ closed, then at least one $F_{n}$ has non-empty interior.
\end{crll} \

\begin{pf}[source=Primary Source Material]
	if $M$ is as above, then $\eset=\bigcap_{n}F_{n}^{c}$.
	exercise: $F_{n}^{c}$ is dense iff $F_{n}$ has no interior.
\end{pf}
we now prove bct.

\newpage
\begin{pf}[source=Primary Source Material]
	let $x\in M$, $\ep>0$. wts $\ex p\in G\cap B(x,\ep)$.
	we'll construct a cauchy seq inductively, w pts $p_{n}$ and radii $r_{n}$ s.t.
	$B(p_{1},2r_{1})\subseteq B(x,\ep)$, and:
	\begin{equation*}
		B(p_{n},2r_{n})\subseteq B(p_{n-1},r_{n-1})\cap \bigcap_{k=1}^{n-1}G_{k}
		\ =: \ B
	\end{equation*}
	to start, pick $p_{1}=x$ and $r_{1}=\ep/2$.
	now, for any $n$, $\bigcap_{k=1}^{n-1}G_{k}$ open dense. thus, $\ex y \in B$.
	since $B$ is open, $\ex r>0$ s.t. $B(y,r)\subseteq B$.
	thus, take $p_{n}=y$ and $r_{n}=\min(r/2,1/n)$.
	since $r_{n}\sto0$, $p_{n}$ is cauchy, thus $p_{n}\sto p$.
	since $p_{k}\in B(p_{n},r_{n})$ for all $k\geq n$, we have:
	\begin{equation*}
		p\in\bar{B(p_{n},r_{n})}\subseteq B(p_{n},2r_{n})
	\end{equation*}
	since $p\in B(p_{n},2r_{n})$ for all $n$, then clearly $p\in B(x,\ep)$ and $p\in G_{n}$ $\fa n$
	as needed.
\end{pf}
kinda missed what he said but sth abt ill-behaved fn's.

\begin{prop}
	$\ex$ a thick $A\subseteq C_{0}([a,b])$ s.t. $\fa f\in A$:
	\begin{enumerate}
		\item the left/right derivs dne for any $x\in[a,b]$
		\item $f$ is not monotone on any $I=(c,d)$.
	\end{enumerate}
\end{prop} \

\begin{pf}[source=Primary Source Material]
	define:
	\begin{gather*}
		\Delta_{h}f(x) \ := \ \frac{f(x+h)-f(x)}{h} \\
		R_{n} \ := \ \set{f:\fa x\in[a,b-1/n],\ex h<1/n,\abs{\Delta_{h}f_{x}}>n} \\
		L_{n} \ := \ \set{f:\fa x\in[a+1/n,b],\ex -1/n<h<0,\abs{\Delta_{h}f_{x}}>n} \\
		G_{n} \ := \ \set{f: f \trm{ not monotone on any interval } I\subseteq[a,b] \trm{ of length } 1/n}
	\end{gather*}
	need to check:
	\begin{enumerate}[i)]
		\item $R_{n},L_{n},G_{n}$ open + dense
		\item $\bigcap_{n=1}^{\infty}R_{n}\cap L_{n}\cap G_{n}$ has the required properties
	\end{enumerate}
	second pt is clear; we check the first pt.
	to see $R_{n}$ open, note:
	\begin{equation*}
		\abs{\Delta_{h}f(x)-\Delta_{h}g(x)}\leq\frac{2\norm{f-g}_{\infty}}{h}
	\end{equation*}
	let $f\in R_{n},x\in[a,b-1/n]$. then $\ex h_{x}$ s.t. $\abs{\Delta_{h_{x}}f(x)}>n$.
	since $f$ cts, $\ex T_{x}\ni x$ open intvl and $v_{x}>0$ s.t. $\fa t\in T_{x}$:
	\begin{equation*}
		\abs{\Delta_{h}f(t)}>n+v_{x}
	\end{equation*}
	in particular, by cptness:
	\begin{equation*}
		[a,b-1/n] \ = \ \bigcup_{x\in[a,b-1/n]}T_{x} \ = \ \bigcup_{j=1}^{m}T_{x_{j}}
	\end{equation*}
	let $h=\min(h_{x_{j}})$ and $v=\min(v_{x_{j}})$.
	then, if $\ep<vh/2$ and $\norm{g-f}_{\infty}<\ep$, we have:
	\begin{equation*}
		\abs{\Delta_{h_{x_{j}}}g(t)} \ > \ \abs{\Delta_{h_{x_{j}}}f(t)}-\frac{2\ep}{h_{x_{j}}}
		\ > \ n+v_{x_{j}}-\frac{vh}{h_{x_{j}}} \ > \ n+v-v \ = \ n
	\end{equation*}
	for $t\in T_{x_{j}}$. in particular, $g\in R_{n}$, so $R_{n}$ and similarly $L_{n}$ open.
	pf cont below.
\end{pf}

\lecdate{(skipped) Lec 13 - Feb 26 (Week 7)}

we now check that $G_{n}$ is open, or $G_{n}^{c}$ closed.
let $f_{m}\in G_{n}^{c}$ s.t. $f_{m}\sto f$.
then $\ex I_{m}=[c_{m},d_{m}]$ with $d_{m}-c_{m}=1/n$ s.t. $f_{m}$ is monotone
on $I_{m}$.
by taking a subseq, sps wlog $c_{m}\sto c$ and $d_{m}\sto d$, and all
$f_{m}$ either increasing or decreasing.
then since $f_{m}\sto f$ ptwise, $f$ monotone on $[a,b]$, so $f\in G_{n}^{c}$.
thus, $G_{n}^{c}$ closed, so $G_{n}$ open.

it remains to see that $R_{n}\cap L_{n}\cap G_{n}$ dense.
by weierstrass approx, it suffices to show that poly's are in the closure.
indeed, let $f$ be a poly.
let $\theta_{m}:\bR\gto\bR$ be the fn s.t. $\theta_{m}(k/m)$ is 0 for even $k$ and 1 for odd $k$,
and linear btwn them (draw a picture).
let $\ep>0$, $f_{m}=f+\ep\theta_{m}/2$.
then $\norm{f-f_{m}}_{\infty}<\ep$, so it suffices to check $f_{m}\in R_{n}\cap L_{n}\cap G_{n}$
for suff. large $m$. for any $x,h$, we have:
\begin{equation*}
	\abs{\Delta_{h}f_{m}(x)}
	\ \geq \ \frac{\ep}{2}\abs{\Delta_{h}\theta_{m}(x)} - \sup_{x\in[a,b]}\abs{f(x)}a
	\ \geq \ \frac{\ep}{2}\abs{\Delta_{h}\theta_{m}(x)}-C
\end{equation*}
for some $C>0$ indep of $h,x$.
then for $x\in\bR$, note (or check) $\ex0<h<2/m$ s.t.:
\begin{equation*}
	\abs{\Delta_{h}\theta_{m}(x)}>\frac{m}{4}
\end{equation*}
thus for suff. large $m$, $f_{m}\in R_{n}$ and similarly $L_{n}$.
to see that $f_{m}\in G_{n}$ for large $m$, note that for any interval $I$ of len $1/n$,
if $m>10n$, $\ex x<y<z$ s.t. $\abs{x-z}<1/m$ and:
\begin{equation*}
	\theta_{m}(y)\geq\theta_{m}(x)+1 \qquad \theta_{m}(z)\leq\theta_{m}(y)-1
\end{equation*}
OTOH, $\abs{f(y)-f(x)}\leq\norm{f'}_{\infty}/m$ and $\abs{f(z)-f(y)}\leq\norm{f'}_{\infty}/m$,
so for suff. large $m$:
\begin{equation*}
	f_{m}(y)\geq f_{m}(x)+\frac{\ep}{4} \qquad f_{m}(z)\leq f_{m}(y)-\frac{\ep}{4}
\end{equation*}
and the result follows.
